% Auto-generated by scripts/python/_outputs/power_calc.py. Do not edit.
\begin{tabular}{lll}
\hline\hline
Hypothesis & Assumption & Minimum detectable effect \\
\hline
H1 (delegation, two-proportion test) & Baseline $p_0=0.30$ & $\Delta = 0.216$ (21.6 pp) \\
H1 (delegation, two-proportion test) & Baseline $p_0=0.50$ & $\Delta = 0.215$ (21.5 pp) \\
H1 (delegation, two-proportion test) & Baseline $p_0=0.70$ & $\Delta = 0.179$ (17.9 pp) \\
\hline
H2 (paired, within-subject) & $\sigma_d = 0.5\,\delta_a = \pounds 0.10$ & $\Delta = \pounds 0.031$ \\
H2 (paired, within-subject) & $\sigma_d = 1\,\delta_a = \pounds 0.20$ & $\Delta = \pounds 0.063$ \\
H2 (paired, within-subject) & $\sigma_d = 2\,\delta_a = \pounds 0.40$ & $\Delta = \pounds 0.125$ \\
H2 (paired, within-subject) & $\sigma_d = 3\,\delta_a = \pounds 0.60$ & $\Delta = \pounds 0.188$ \\
H2 (paired, within-subject) & $\sigma_d = 4\,\delta_a = \pounds 0.80$ & $\Delta = \pounds 0.251$ \\
\hline\hline
\multicolumn{3}{p{0.95\textwidth}}{\footnotesize \textit{Note:} Both tests assume $n = 80$ per role per condition, $\alpha = 0.05$ two-sided, power $= 0.80$. For H2, $\sigma_d$ is the within-subject standard deviation of the punishment difference (delegated minus non-delegated, conditional on bad outcome) and is expressed as a multiple of the anchor effect $\delta_a = \pounds 0.20$ (the equivalence bound used in Result~2 in the body of the paper, equal to $10\%$ of the maximum punishment of $\pounds 2$). $\Delta$ is the corresponding minimum detectable effect in pounds.}
\end{tabular}
